% Every way TeX will read an integer. See tools/battery.sh.
\input plain
\catcode`\@=11 \catcode`\:=11 \catcode`\_=11
\count5=17 \dimen5=3pt
\def\show#1{\count0=#1 \message{<\the\count0>}}
\show{123}
\show{-42}
\show{--5}
\show{+ 7}
\show{`A}
\show{`\A}
\show{"FF}
\show{'77}
\show{\count5}
\show{-\count5}
\show{\numexpr 2*3\relax}
\show{\numexpr 7/2\relax}
\show{\catcode`\A}
\message{<\the\numexpr 123\relax>}
\message{<\the\numexpr `A\relax>}
\message{<\the\numexpr "FF\relax>}
\message{<\the\numexpr \count5\relax>}
\message{<\the\numexpr 2*3+1\relax>}
% Number scanning is execution, not lexing: TeX reads digits with get_x_token
% (tex.web 445), so a macro standing for a digit is absorbed into the number.
% expl3's whole f-type expansion rests on this, written `\romannumeral-`0`:
% the value is negative so nothing is emitted, but scanning it expands what
% follows until something unexpandable stops the scan.
\def\z{5}\def\foo{ABC}\def\pair{\z 7}
\message{<\romannumeral1\z>}
\message{<\romannumeral-`0\foo>}
\message{<\number1\z>}
\message{<\number3\pair>}
\message{<\the\numexpr 1\z\relax>}
% Signs bind to every form of number, not only to digits -- which is why the
% sign scan has to come first.
\count5=17
\message{<\number-`A>}
\message{<\number-"FF>}
\message{<\number-'77>}
\message{<\number-\count5>}
% A \chardef'd name is an internal integer, and expl3's controlled expansion
% relies on it: \exp:w is \romannumeral and \exp_end: is a chardef'd zero, so
% `\exp:w <tokens> ` expands <tokens> while scanning a number and terminates
% when one of them leaves \exp_end: in front.
\chardef\c_zero_int=0
\let\exp:w\romannumeral \let\exp_end:\c_zero_int
\def\lvl@one{\exp_end: ABC}\def\lvl@two{\lvl@three}\def\lvl@three{\exp_end: XY}
\message{<\exp:w \lvl@one>}
\message{<\exp:w \lvl@two>}
\message{<\number\c_zero_int>}
\chardef\half=50
\message{<\number\half>}
\message{<\romannumeral\half>}
\message{<\the\numexpr\half+1\relax>}
\bye
